\documentclass[11pt,letterpaper]{article}
\usepackage[utf8]{inputenc}
\usepackage[margin=0.75in]{geometry}
\usepackage{amsmath,amssymb,amsthm}
\usepackage{graphicx}
\usepackage{xcolor}
\usepackage{enumitem}
\usepackage{booktabs}
\usepackage{tabularx}
\usepackage{longtable}
\usepackage{multicol}
\usepackage{tcolorbox}
\usepackage{fancyhdr}
\usepackage{titlesec}
\usepackage[hidelinks]{hyperref}

% Colors
\definecolor{criticalcolor}{RGB}{220,20,60}
\definecolor{highcolor}{RGB}{255,140,0}
\definecolor{mediumcolor}{RGB}{70,130,180}
\definecolor{headercolor}{RGB}{25,25,112}

% Header and Footer
\pagestyle{fancy}
\fancyhf{}
\fancyhead[L]{\small ORF309 Midterm 2 Study Guide}
\fancyhead[R]{\small March 31, 2026}
\fancyfoot[C]{\thepage}

% Section formatting
\titleformat{\section}{\Large\bfseries\color{headercolor}}{\thesection}{1em}{}
\titleformat{\subsection}{\large\bfseries\color{headercolor}}{\thesubsection}{1em}{}

% Custom boxes
\newtcolorbox{criticalbox}{
  colback=criticalcolor!5,
  colframe=criticalcolor,
  fonttitle=\bfseries,
  title=CRITICAL PRIORITY
}

\newtcolorbox{highbox}{
  colback=highcolor!5,
  colframe=highcolor,
  fonttitle=\bfseries,
  title=HIGH PRIORITY
}

\newtcolorbox{tipbox}{
  colback=blue!5,
  colframe=blue!70,
  fonttitle=\bfseries,
  title=Key Insight
}

\begin{document}

% Title Page
\begin{titlepage}
\centering
\vspace*{2cm}

{\Huge\bfseries ORF309 Midterm 2\par}
\vspace{0.5cm}
{\LARGE\bfseries Comprehensive Study Guide\par}
\vspace{2cm}

{\Large Analysis of 7 Past Exams\par}
\vspace{0.3cm}
{\large Topics, Frequencies, and Learning Strategy\par}

\vspace{3cm}

{\large Created: March 31, 2026\par}

\vfill

{\large\textbf{Exam Coverage:} Slides 7--23 (Post-Midterm 1)\par}
\vspace{0.5cm}
{\large\textbf{Exams Analyzed:} F2025, S2025, F2022, S2022, 2021, 2019, 2017\par}

\end{titlepage}

\tableofcontents
\newpage

\section{Executive Summary}

\begin{tipbox}
\textbf{Bottom Line:} This exam is highly predictable! Three topics appear in 100\% of exams and account for $\sim$75\% of the total points. Master these three techniques and you'll be well-prepared.
\end{tipbox}

\subsection{Critical Statistics}

Based on analysis of 7 past Midterm 2 exams:

\begin{itemize}[leftmargin=*]
  \item \textbf{Slides Covered:} 7--23 (everything after Midterm 1)
  \item \textbf{Exam Format:} 3 questions, 30 points each, 90 total points
  \item \textbf{Time:} 105 minutes for problems + 15 minutes for submission
  \item \textbf{Most Predictable Topics:} Poisson Processes, Random Walks, Joint Densities
\end{itemize}

\subsection{The Big Three Topics}

\begin{enumerate}[leftmargin=*]
  \item \textbf{Poisson Processes} (Slides 16--19): Appears in \textcolor{criticalcolor}{7/7 exams (100\%)}
  \begin{itemize}
    \item Average weight: $\sim$40\% of exam
    \item Key techniques: Thinning, superposition, competing processes, conditional expectations
  \end{itemize}

  \item \textbf{Random Walks \& First-Step Analysis} (Slides 20--22): Appears in \textcolor{criticalcolor}{7/7 exams (100\%)}
  \begin{itemize}
    \item Average weight: $\sim$35\% of exam
    \item Key techniques: Hitting probabilities, expected hitting times, gambler's ruin
  \end{itemize}

  \item \textbf{Joint \& Conditional Densities} (Slides 11--13): Appears in \textcolor{criticalcolor}{6/7 exams (86\%)}
  \begin{itemize}
    \item Average weight: $\sim$25\% of exam
    \item Key techniques: Building joint PDFs, conditional expectations, law of total expectation
  \end{itemize}
\end{enumerate}

\subsection{Also Important}

\begin{itemize}[leftmargin=*]
  \item \textbf{Exponential Distribution \& Lifetimes} (Slides 11, 14--15): 5/7 exams (71\%)
  \item \textbf{Brownian Motion} (Slide 23): 3/7 exams (43\%), increasingly common in recent exams
\end{itemize}

\newpage

\section{Topic Frequency Analysis}

\begin{table}[h]
\centering
\caption{Topic Appearance Frequency Across 7 Past Exams}
\begin{tabularx}{\textwidth}{Xlcc}
\toprule
\textbf{Topic} & \textbf{Frequency} & \textbf{Slides} & \textbf{Priority} \\
\midrule
Poisson Processes & \textcolor{criticalcolor}{7/7 (100\%)} & 16--19 & CRITICAL \\
Random Walks & \textcolor{criticalcolor}{7/7 (100\%)} & 20--22 & CRITICAL \\
Joint/Conditional Densities & \textcolor{criticalcolor}{6/7 (86\%)} & 11--13 & CRITICAL \\
Law of Total Expectation & \textcolor{criticalcolor}{6/7 (86\%)} & 11--12 & CRITICAL \\
Thinning Property & \textcolor{criticalcolor}{6/7 (86\%)} & 17 & CRITICAL \\
First-Step Analysis & \textcolor{criticalcolor}{7/7 (100\%)} & 20--21 & CRITICAL \\
Exponential/Lifetimes & \textcolor{highcolor}{5/7 (71\%)} & 11, 14--15 & HIGH \\
Competing Processes & \textcolor{highcolor}{5/7 (71\%)} & 16--17 & HIGH \\
Biased Random Walks & \textcolor{highcolor}{4/7 (57\%)} & 22 & HIGH \\
Brownian Motion & \textcolor{mediumcolor}{3/7 (43\%)} & 23 & MEDIUM-HIGH \\
\bottomrule
\end{tabularx}
\end{table}

\section{Critical Priority Topics}

\begin{criticalbox}
Learn these topics FIRST. They appear in 80\%+ of exams and account for $\sim$75\% of exam content.
\end{criticalbox}

\subsection{Poisson Processes}

\textbf{Frequency:} 7/7 exams (100\%) \quad \textbf{Weight:} $\sim$40\% of exam \quad \textbf{Slides:} 16--19

\subsubsection{Key Concepts}

\paragraph{1. Basic Poisson Definition}
If arrivals happen at rate $\lambda$ (per unit time):
\begin{itemize}
  \item Number of arrivals $N(t) \sim \text{Poisson}(\lambda t)$
  \item $P(N(t) = k) = \frac{(\lambda t)^k e^{-\lambda t}}{k!}$
  \item $E[N(t)] = \lambda t$, \quad $\text{Var}(N(t)) = \lambda t$
\end{itemize}

\paragraph{2. Thinning Property} \textcolor{criticalcolor}{(APPEARS IN 6/7 EXAMS)}

\begin{tcolorbox}[colback=yellow!10,colframe=orange!80,title=Thinning - Most Important Concept]
If customers arrive at rate $\lambda$ and each customer is ``type A'' with probability $p$ (independently):
\[
\text{Type A arrivals form Poisson process with rate } \lambda p
\]

\textbf{Example:} Customers arrive at rate 2/min, 40\% order coffee
\[
\Rightarrow \text{Coffee orders are Poisson}(2 \times 0.4 = 0.8\text{/min})
\]
\end{tcolorbox}

\paragraph{3. Superposition/Merging}
Multiple independent Poisson processes combine:
\[
\text{If } N_1(t) \sim \text{Poisson}(\lambda_1 t), \, N_2(t) \sim \text{Poisson}(\lambda_2 t) \text{ independent}
\]
\[
\Rightarrow N_1(t) + N_2(t) \sim \text{Poisson}((\lambda_1 + \lambda_2)t)
\]

\paragraph{4. Competing Processes} \textcolor{highcolor}{(APPEARS IN 5/7 EXAMS)}

Process A at rate $\lambda_A$, Process B at rate $\lambda_B$ (independent):
\begin{align*}
P(\text{A happens first}) &= \frac{\lambda_A}{\lambda_A + \lambda_B} \\
E[\text{time until first event}] &= \frac{1}{\lambda_A + \lambda_B}
\end{align*}

\paragraph{5. Conditional Expectations} \textcolor{criticalcolor}{(APPEARS IN 5/7 EXAMS)}

\begin{tcolorbox}[colback=yellow!10,colframe=orange!80,title=Conditional Expectation with Poisson]
Given that exactly $n$ events occurred:
\[
E[\text{Time} \mid N(T) = n] = \frac{n}{\lambda}
\]

\textbf{Example:} Given 10 customers arrived at rate 2/min:
\[
E[T \mid N(T) = 10] = \frac{10}{2} = 5 \text{ minutes}
\]
\end{tcolorbox}

\subsubsection{Example Question Types}

\begin{itemize}[leftmargin=*]
  \item \textbf{F2025 Q3:} Coffee shop with multiple product types (thinning + conditional expectation)
  \item \textbf{S2025 Q3:} Taxi arrivals vs passenger arrivals (competing processes)
  \item \textbf{2022F Q1:} Turkey arrivals with weight thresholds (thinning with exponential)
  \item \textbf{2022S Q1:} Chicken crossing road (conditional expectation)
  \item \textbf{2019 Q1:} Plagues arriving (thinning with acceptance probability)
\end{itemize}

\newpage

\subsection{Random Walks \& First-Step Analysis}

\textbf{Frequency:} 7/7 exams (100\%) \quad \textbf{Weight:} $\sim$35\% of exam \quad \textbf{Slides:} 20--22

\subsubsection{Key Concepts}

\paragraph{Random Walk Setup}
\begin{itemize}
  \item Start at position $S_0$ (usually between barriers $a$ and $b$)
  \item Each step: move $+1$, $-1$, or stay, with given probabilities
  \item Continue until hitting $a$ or $b$ (boundaries)
  \item Hitting time: $T = \min\{n \geq 0 : S_n = a \text{ or } S_n = b\}$
\end{itemize}

\paragraph{First-Step Analysis for HITTING PROBABILITY}

\begin{tcolorbox}[colback=yellow!10,colframe=orange!80,title=First-Step Analysis - Hitting Probability]
Define: $p(i) = P(\text{hit } b \text{ before } a \mid \text{start at } i)$

\textbf{Recursive formula:}
\[
p(i) = p_{\text{up}} \cdot p(i+1) + p_{\text{down}} \cdot p(i-1) + p_{\text{stay}} \cdot p(i)
\]

\textbf{Boundary conditions:}
\[
p(a) = 0, \quad p(b) = 1
\]

\textbf{Solve the system of equations for all intermediate states!}
\end{tcolorbox}

\paragraph{First-Step Analysis for EXPECTED TIME}

\begin{tcolorbox}[colback=red!10,colframe=red!80,title=First-Step Analysis - Expected Time]
Define: $\mu(i) = E[\text{time to hit boundary} \mid \text{start at } i]$

\textbf{Recursive formula:}
\[
\mu(i) = \mathbf{1} + p_{\text{up}} \cdot \mu(i+1) + p_{\text{down}} \cdot \mu(i-1) + p_{\text{stay}} \cdot \mu(i)
\]

\textbf{DON'T FORGET THE ``+1''!} (Most common mistake)

\textbf{Boundary conditions:}
\[
\mu(a) = 0, \quad \mu(b) = 0
\]
\end{tcolorbox}

\paragraph{Special Case: Symmetric Random Walk}

When $P(+1) = P(-1) = 1/2$:
\begin{align*}
p(i) &= \frac{i - a}{b - a} \\
\mu(i) &= (i - a)(b - i)
\end{align*}

\paragraph{Biased Random Walk} \textcolor{highcolor}{(APPEARS IN 4/7 EXAMS)}

When $P(+1) \neq P(-1)$:
\begin{itemize}
  \item Use first-step analysis and solve system of equations
  \item OR use characteristic equation: $r = pr + q/r$ where $p = P(+1)$, $q = P(-1)$
  \item Solve: $pr^2 - r + q = 0 \Rightarrow r = \frac{q}{p}$ (if $p \neq q$)
\end{itemize}

\subsubsection{Problem-Solving Steps}

\begin{enumerate}[leftmargin=*]
  \item Draw a diagram with states $a$, $i$, $b$
  \item Identify what you're solving for: $p(i)$ or $\mu(i)$
  \item Write the recursive formula (don't forget ``+1'' for $\mu$!)
  \item Write boundary conditions
  \item Solve the system of equations (usually 2--3 equations)
\end{enumerate}

\subsubsection{Example Question Types}

\begin{itemize}[leftmargin=*]
  \item \textbf{F2025 Q2:} Gambler with laziness parameter (modified random walk)
  \item \textbf{S2025 Q1:} Flight price movement (biased walk with first-step analysis)
  \item \textbf{2022F Q3:} Indigestion problem (hitting time analysis)
  \item \textbf{2022S Q2:} Easter bunny with $\pm 1$, $\pm 2$ jumps (modified random walk)
  \item \textbf{2021 Q3:} Time warp dance (asymmetric steps)
  \item \textbf{2019 Q3:} Israelites in desert (biased walk with drift)
  \item \textbf{2017 Q3:} Asymmetric casino game (biased walk with $+1/-2$)
\end{itemize}

\newpage

\subsection{Continuous Random Variables: Joint \& Conditional Densities}

\textbf{Frequency:} 6/7 exams (86\%) \quad \textbf{Weight:} $\sim$25\% of exam \quad \textbf{Slides:} 11--13

\subsubsection{Key Concepts}

\paragraph{1. Joint PDF $f_{X,Y}(x,y)$}
\begin{itemize}
  \item Probability density for two continuous random variables
  \item Must integrate to 1: $\int_{-\infty}^{\infty} \int_{-\infty}^{\infty} f_{X,Y}(x,y) \, dx \, dy = 1$
  \item \textbf{Always specify the range!} (e.g., ``for $0 < y < x < \infty$'')
\end{itemize}

\paragraph{2. Marginal PDF from Joint}
\begin{align*}
f_X(x) &= \int_{-\infty}^{\infty} f_{X,Y}(x,y) \, dy \quad \text{(integrate over } y\text{)} \\
f_Y(y) &= \int_{-\infty}^{\infty} f_{X,Y}(x,y) \, dx \quad \text{(integrate over } x\text{)}
\end{align*}

\paragraph{3. Conditional PDF}
\begin{align*}
f_{X|Y}(x|y) &= \frac{f_{X,Y}(x,y)}{f_Y(y)} \\
f_{Y|X}(y|x) &= \frac{f_{X,Y}(x,y)}{f_X(x)}
\end{align*}

\paragraph{4. Building Joint PDF from Conditional}
\begin{align*}
f_{X,Y}(x,y) &= f_{X|Y}(x|y) \cdot f_Y(y) \\
f_{X,Y}(x,y) &= f_{Y|X}(y|x) \cdot f_X(x)
\end{align*}

\paragraph{5. Law of Total Expectation} \textcolor{criticalcolor}{(SUPER IMPORTANT)}

\begin{tcolorbox}[colback=yellow!10,colframe=orange!80,title=Law of Total Expectation (Tower Property)]
\[
E[Y] = E[E[Y|X]]
\]

In integral form:
\[
E[Y] = \int_{-\infty}^{\infty} E[Y|X=x] \cdot f_X(x) \, dx
\]

\textbf{Incredibly useful for computing expectations!}
\end{tcolorbox}

\paragraph{6. Common Transformation: $Y = UX$}

When $U \sim \text{Uniform}[0,1]$ independent of $X$:
\begin{itemize}
  \item Given $X = x$, $Y$ is uniform on $[0, x]$
  \item $f_{Y|X}(y|x) = \frac{1}{x}$ for $0 < y < x$
  \item $f_{X,Y}(x,y) = f_{Y|X}(y|x) \cdot f_X(x) = \frac{1}{x} f_X(x)$ for $0 < y < x$
\end{itemize}

\subsubsection{Problem-Solving Strategy}

\begin{enumerate}[leftmargin=*]
  \item Identify what's given (conditional or marginal?)
  \item Build joint PDF using $f_{X,Y}(x,y) = f_{X|Y}(x|y) \cdot f_Y(y)$ or $f_{Y|X}(y|x) \cdot f_X(x)$
  \item \textbf{Check the range carefully!} Draw a picture if needed
  \item For expectations: use law of total expectation $E[Y] = E[E[Y|X]]$
\end{enumerate}

\subsubsection{Example Question Types}

\begin{itemize}[leftmargin=*]
  \item \textbf{F2025 Q1:} Insurance payments $Y = UX$ (joint density, conditional probability)
  \item \textbf{S2025 Q2:} Flight delay with weather index (joint density, conditional probabilities)
  \item \textbf{2022F Q2:} Pie eaten by dog (conditional expectation with gamma distribution)
\end{itemize}

\newpage

\section{High Priority Topics}

\begin{highbox}
Learn these topics after mastering the Critical Priority topics. They appear in 40--80\% of exams.
\end{highbox}

\subsection{Exponential Distribution \& Lifetimes}

\textbf{Frequency:} 5/7 exams (71\%) \quad \textbf{Weight:} $\sim$15\% of exam \quad \textbf{Slides:} 11, 14--15

\subsubsection{Exponential Distribution Basics}

If $T \sim \text{Exp}(\lambda)$:
\begin{align*}
\text{PDF:} \quad & f(t) = \lambda e^{-\lambda t} \quad \text{for } t > 0 \\
\text{CDF:} \quad & F(t) = 1 - e^{-\lambda t} \\
\text{Survival:} \quad & P(T > t) = e^{-\lambda t} \\
\text{Mean:} \quad & E[T] = \frac{1}{\lambda} \\
\text{Variance:} \quad & \text{Var}(T) = \frac{1}{\lambda^2}
\end{align*}

\subsubsection{Memoryless Property}

\begin{tcolorbox}[colback=yellow!10,colframe=orange!80,title=Memoryless Property - KEY!]
\[
P(T > s + t \mid T > s) = P(T > t)
\]

\textbf{Interpretation:} The distribution ``doesn't remember the past''

\textbf{Key fact:} Exponential is the ONLY continuous distribution with this property!
\end{tcolorbox}

\subsubsection{Minimum of Independent Exponentials}

If $T_1 \sim \text{Exp}(\lambda_1)$ and $T_2 \sim \text{Exp}(\lambda_2)$ are independent:
\begin{align*}
\min(T_1, T_2) &\sim \text{Exp}(\lambda_1 + \lambda_2) \\
P(T_1 < T_2) &= \frac{\lambda_1}{\lambda_1 + \lambda_2} \\
E[\min(T_1, T_2)] &= \frac{1}{\lambda_1 + \lambda_2}
\end{align*}

\textbf{Connection to Poisson:} This relates to competing Poisson processes!

\subsubsection{Maximum of Exponentials}

\begin{align*}
E[\max(T_1, T_2)] &= \frac{1}{\lambda_1} + \frac{1}{\lambda_2} - \frac{1}{\lambda_1 + \lambda_2}
\end{align*}

No simple closed form for the distribution of $\max(T_1, T_2)$.

\subsubsection{Example Question Types}

\begin{itemize}[leftmargin=*]
  \item \textbf{2022S Q1:} Chicken crossing with exponential wait time
  \item \textbf{2019 Q2:} Bread rising and proofing (max of exponentials)
  \item \textbf{2017 Q2:} Smoker vs non-smoker lifetimes (hazard rates)
\end{itemize}

\subsection{Brownian Motion}

\textbf{Frequency:} 3/7 exams (43\%) \quad \textbf{Weight:} $\sim$20\% when it appears \quad \textbf{Slides:} 23

\textbf{Note:} Increasingly common in recent exams!

\subsubsection{Standard Brownian Motion Definition}

$B(t)$ is a \textit{standard Brownian motion} if:
\begin{enumerate}
  \item $B(0) = 0$ (starts at origin)
  \item $B(t)$ has continuous paths
  \item $B(t)$ has independent increments
  \item $B(t) - B(s) \sim N(0, t-s)$ for $t > s$
\end{enumerate}

\subsubsection{Key Properties}

\begin{align*}
E[B(t)] &= 0 \\
\text{Var}(B(t)) &= t \\
B(t) &\sim N(0, t) \\
E[B(t)^2] &= t \quad \text{(since } \text{Var} = E[X^2] - (E[X])^2 = E[X^2]\text{)}
\end{align*}

\subsubsection{Conditional Expectations}

For $t > s$:
\begin{align*}
E[B(t) \mid B(s) = x] &= x \\
E[B(t)^2 \mid B(s) = x] &= x^2 + (t - s)
\end{align*}

\subsubsection{Common Application: Stock Prices}

Geometric Brownian motion:
\[
S(t) = S_0 \cdot e^{B(t)}
\]
Used extensively in financial modeling.

\subsubsection{Example Question Types}

\begin{itemize}[leftmargin=*]
  \item \textbf{2022S Q3:} Stock price $S(t) = e^{B(t)}$ (probabilities with BM)
  \item \textbf{2021 Q1:} Two ants doing independent Brownian motion (distance calculation)
\end{itemize}

\newpage

\section{Key Formulas \& Techniques}

\subsection{Poisson Processes}

\begin{tcolorbox}[colback=blue!5,colframe=blue!70,title=Poisson Process Formulas]
If $N(t) \sim \text{Poisson}(\lambda t)$:
\begin{align*}
P(N(t) = k) &= \frac{(\lambda t)^k e^{-\lambda t}}{k!} \\
E[N(t)] &= \lambda t \\
\text{Var}(N(t)) &= \lambda t
\end{align*}

\textbf{Thinning:} If arrivals are rate $\lambda$, each type A with probability $p$:
\[
\Rightarrow \text{Type A arrivals are Poisson}(\lambda p)
\]

\textbf{Competing:}
\[
P(\text{A before B}) = \frac{\lambda_A}{\lambda_A + \lambda_B}
\]

\textbf{Conditional:}
\[
E[T \mid N(T) = n] = \frac{n}{\lambda}
\]
\end{tcolorbox}

\subsection{Random Walk First-Step Analysis}

\begin{tcolorbox}[colback=blue!5,colframe=blue!70,title=First-Step Analysis Formulas]
\textbf{For hitting probability $p(i)$:}
\[
p(i) = p_{\text{up}} \cdot p(i+1) + p_{\text{down}} \cdot p(i-1) + p_{\text{stay}} \cdot p(i)
\]
\textbf{Boundaries:} $p(a) = 0$, $p(b) = 1$

\vspace{0.3cm}

\textbf{For expected hitting time $\mu(i)$:}
\[
\mu(i) = \mathbf{1} + p_{\text{up}} \cdot \mu(i+1) + p_{\text{down}} \cdot \mu(i-1) + p_{\text{stay}} \cdot \mu(i)
\]
\textbf{Boundaries:} $\mu(a) = 0$, $\mu(b) = 0$

\vspace{0.3cm}

\textbf{For symmetric walk ($p = q = 1/2$):}
\begin{align*}
p(i) &= \frac{i - a}{b - a} \\
\mu(i) &= (i - a)(b - i)
\end{align*}
\end{tcolorbox}

\subsection{Exponential Distribution}

\begin{tcolorbox}[colback=blue!5,colframe=blue!70,title=Exponential Distribution Formulas]
If $T \sim \text{Exp}(\lambda)$:
\begin{align*}
f(t) &= \lambda e^{-\lambda t} \\
E[T] &= \frac{1}{\lambda} \\
P(T > t) &= e^{-\lambda t}
\end{align*}

\textbf{Memoryless:} $P(T > s+t \mid T > s) = P(T > t)$

\vspace{0.3cm}

\textbf{Min of exponentials:}
\begin{align*}
\min(T_1, T_2) &\sim \text{Exp}(\lambda_1 + \lambda_2) \\
P(T_1 < T_2) &= \frac{\lambda_1}{\lambda_1 + \lambda_2}
\end{align*}
\end{tcolorbox}

\subsection{Continuous Random Variables}

\begin{tcolorbox}[colback=blue!5,colframe=blue!70,title=Joint \& Conditional Density Formulas]
\textbf{Joint to Marginal:}
\[
f_X(x) = \int_{-\infty}^{\infty} f_{X,Y}(x,y) \, dy
\]

\textbf{Conditional:}
\[
f_{X|Y}(x|y) = \frac{f_{X,Y}(x,y)}{f_Y(y)}
\]

\textbf{Law of Total Expectation:}
\begin{align*}
E[Y] &= E[E[Y|X]] \\
&= \int_{-\infty}^{\infty} E[Y|X=x] \cdot f_X(x) \, dx
\end{align*}
\end{tcolorbox}

\subsection{Brownian Motion}

\begin{tcolorbox}[colback=blue!5,colframe=blue!70,title=Brownian Motion Formulas]
If $B(t)$ is standard Brownian motion:
\begin{align*}
B(0) &= 0 \\
E[B(t)] &= 0 \\
\text{Var}(B(t)) &= t \\
B(t) &\sim N(0,t) \\
B(t) - B(s) &\sim N(0, t-s) \quad \text{for } t > s
\end{align*}

\textbf{Conditional:}
\begin{align*}
E[B(t)^2 \mid B(s) = x] &= x^2 + (t-s) \quad \text{for } t > s
\end{align*}
\end{tcolorbox}

\newpage

\section{Common Mistakes to Avoid}

\subsection{Poisson Processes}

\begin{itemize}[leftmargin=*]
  \item[\textcolor{red}{$\times$}] Forgetting to use thinning when events split into types
  \item[\textcolor{red}{$\times$}] Not recognizing when to use conditional expectation $E[T \mid N = n]$
  \item[\textcolor{green!60!black}{$\checkmark$}] Always check if arrivals split by type $\rightarrow$ use thinning
\end{itemize}

\subsection{Random Walks}

\begin{itemize}[leftmargin=*]
  \item[\textcolor{red}{$\times$}] Forgetting the ``+1'' in the expected time formula: $\mu(i) = \mathbf{1} + \cdots$
  \item[\textcolor{red}{$\times$}] Wrong boundary conditions (should be 0 at both barriers for $\mu$)
  \item[\textcolor{green!60!black}{$\checkmark$}] Write out the recursive formula carefully with all probabilities
\end{itemize}

\subsection{Joint Densities}

\begin{itemize}[leftmargin=*]
  \item[\textcolor{red}{$\times$}] Not checking the range/support of the distribution
  \item[\textcolor{red}{$\times$}] Forgetting to normalize (density must integrate to 1)
  \item[\textcolor{green!60!black}{$\checkmark$}] Always specify the region where $f(x,y) > 0$
\end{itemize}

\subsection{Exponential}

\begin{itemize}[leftmargin=*]
  \item[\textcolor{red}{$\times$}] Confusing rate $\lambda$ with mean (mean = $1/\lambda$)
  \item[\textcolor{red}{$\times$}] Not using memoryless property when applicable
  \item[\textcolor{green!60!black}{$\checkmark$}] Remember: min of exponentials $\rightarrow$ sum of rates
\end{itemize}

\newpage

\section{Detailed Exam Breakdown}

\subsection{Fall 2025 Midterm}

\textbf{Q1 (Insurance --- 30 pts):} Joint/Conditional Densities, Transformations ($Y = UX$)
\begin{itemize}
  \item Topics: Continuous RVs, conditional probability, tail distributions
  \item Slides: 11--13
\end{itemize}

\textbf{Q2 (Sleepy Gambler --- 30 pts):} Modified Random Walk with laziness parameter
\begin{itemize}
  \item Topics: First-step analysis, hitting probabilities, expected hitting times
  \item Slides: 20--21
\end{itemize}

\textbf{Q3 (Coffee Shop --- 30 pts):} Poisson with Thinning, Conditional Expectation
\begin{itemize}
  \item Topics: Thinning property, conditional expectations, Poisson processes
  \item Slides: 16--18
\end{itemize}

\subsection{Spring 2025 Midterm}

\textbf{Q1 (Flight Prices --- 30 pts):} Biased Random Walk with barriers
\begin{itemize}
  \item Topics: First-step analysis, hitting probabilities, expected times
  \item Slides: 20--22
\end{itemize}

\textbf{Q2 (Weather Index --- 30 pts):} Joint Density, Exponential conditioning
\begin{itemize}
  \item Topics: Joint PDFs, conditional probabilities with continuous RVs
  \item Slides: 11--13, 14--15
\end{itemize}

\textbf{Q3 (Taxi Stop --- 30 pts):} Competing Poisson Processes
\begin{itemize}
  \item Topics: Poisson processes, conditional expectations, competing processes
  \item Slides: 16--18
\end{itemize}

\subsection{Fall 2022 Midterm}

\textbf{Q1 (Turkeys --- 30 pts):} Poisson with Thinning, Competing Processes
\begin{itemize}
  \item Topics: Thinning (exponential threshold), competing Poisson vs customer arrivals
  \item Slides: 16--18
\end{itemize}

\textbf{Q2 (Pie --- 30 pts):} Conditional Expectation, Gamma Distribution
\begin{itemize}
  \item Topics: Law of total expectation, conditional expectations with continuous RVs
  \item Slides: 11--13
\end{itemize}

\textbf{Q3 (Indigestion --- 30 pts):} Random Walk, Hitting Time
\begin{itemize}
  \item Topics: Discrete random walk, probability of never hitting barrier
  \item Slides: 20--21
\end{itemize}

\subsection{Spring 2022 Midterm}

\textbf{Q1 (Chicken Crossing --- 30 pts):} Poisson + Exponential, Conditional Expectation
\begin{itemize}
  \item Topics: Poisson count, conditional expectation given count
  \item Slides: 16--18, 11
\end{itemize}

\textbf{Q2 (Easter Bunny --- 30 pts):} Random Walk with jumps of $\pm 1$, $\pm 2$
\begin{itemize}
  \item Topics: Modified random walk, first-step analysis
  \item Slides: 20--22
\end{itemize}

\textbf{Q3 (Stock Price --- 30 pts):} Brownian Motion, $e^{B(t)}$
\begin{itemize}
  \item Topics: Brownian motion properties, probabilities with BM
  \item Slides: 23
\end{itemize}

\subsection{2021 Midterm}

\textbf{Q1 (Tiger-ants --- 30 pts):} Two Independent Brownian Motions
\begin{itemize}
  \item Topics: Conditional expectations with BM, distance between BMs
  \item Slides: 23
\end{itemize}

\textbf{Q2 (Fruity Yogurt --- 30 pts):} Poisson with Thinning (two-stage)
\begin{itemize}
  \item Topics: Thinning with multiple probabilities, conditional expectations
  \item Slides: 16--18
\end{itemize}

\textbf{Q3 (Time Warp --- 30 pts):} Modified Random Walk
\begin{itemize}
  \item Topics: First-step analysis with asymmetric steps
  \item Slides: 20--22
\end{itemize}

\subsection{2019 Midterm}

\textbf{Q1 (Plagues --- 30 pts):} Poisson with Thinning/Acceptance
\begin{itemize}
  \item Topics: Thinning with acceptance probability, types of events
  \item Slides: 16--18
\end{itemize}

\textbf{Q2 (Matzot --- 30 pts):} Max of Two Exponentials
\begin{itemize}
  \item Topics: Expected max, probability calculations with exponentials
  \item Slides: 11, 14--15
\end{itemize}

\textbf{Q3 (Desert --- 30 pts):} Random Walk with Drift
\begin{itemize}
  \item Topics: Biased random walk, first-step analysis
  \item Slides: 20--22
\end{itemize}

\subsection{2017 Midterm}

\textbf{Q1 (Twitter --- 30 pts):} Multiple Competing Poisson Processes
\begin{itemize}
  \item Topics: Superposition, competing processes, conditional expectations
  \item Slides: 16--18
\end{itemize}

\textbf{Q2 (Smoking --- 30 pts):} Hazard Functions, Lifetimes
\begin{itemize}
  \item Topics: Hazard rate, probability one outlasts another
  \item Slides: 14--15
\end{itemize}

\textbf{Q3 (Asymmetric Game --- 30 pts):} Biased Random Walk ($+1/-2$)
\begin{itemize}
  \item Topics: First-step analysis with asymmetric steps
  \item Slides: 20--22
\end{itemize}

\newpage

\section{Study Strategy \& Learning Roadmap}

\subsection{Suggested 4-Day Study Plan}

\subsubsection{Day 1: Poisson Processes (6--8 hours)}

\textbf{Morning (3--4 hours):}
\begin{itemize}[leftmargin=*]
  \item Read Slides 16--19 carefully
  \item Take notes on thinning, superposition, competing processes
  \item Work through slide examples
\end{itemize}

\textbf{Afternoon (3--4 hours):}
\begin{itemize}[leftmargin=*]
  \item Practice F2025 Q3 (Coffee shop)
  \item Practice S2025 Q3 (Taxi stop)
  \item Practice 2019 Q1 (Plagues)
  \item Check solutions, understand every step
\end{itemize}

\textbf{Evening Review:}
\begin{itemize}[leftmargin=*]
  \item Can you explain thinning without looking at notes?
  \item Can you set up competing process problems?
\end{itemize}

\subsubsection{Day 2: Random Walks (6--8 hours)}

\textbf{Morning (3--4 hours):}
\begin{itemize}[leftmargin=*]
  \item Read Slides 20--22 carefully
  \item Focus on first-step analysis technique
  \item Write out formula templates
\end{itemize}

\textbf{Afternoon (3--4 hours):}
\begin{itemize}[leftmargin=*]
  \item Practice F2025 Q2 (Sleepy gambler)
  \item Practice S2025 Q1 (Flight prices)
  \item Practice 2017 Q3 (Asymmetric game)
\end{itemize}

\textbf{Evening Review:}
\begin{itemize}[leftmargin=*]
  \item Write first-step formulas from memory
  \item Check: do you remember the ``+1'' for expected time?
\end{itemize}

\subsubsection{Day 3: Joint Densities + Mixed Practice (6--8 hours)}

\textbf{Morning (2--3 hours):}
\begin{itemize}[leftmargin=*]
  \item Read Slides 11--13
  \item Focus on conditional densities and law of total expectation
  \item Work through examples
\end{itemize}

\textbf{Afternoon (3--4 hours):}
\begin{itemize}[leftmargin=*]
  \item Practice F2025 Q1 (Insurance)
  \item Practice S2025 Q2 (Weather index)
  \item Do a FULL timed practice exam (F2025 or S2025)
\end{itemize}

\textbf{Evening (1--2 hours):}
\begin{itemize}[leftmargin=*]
  \item Quick review of Exponential (Slides 11, 14--15)
  \item If time: Brownian Motion basics (Slide 23)
\end{itemize}

\subsubsection{Day 4: Final Review + Weak Areas (4--6 hours)}

\textbf{Morning (2--3 hours):}
\begin{itemize}[leftmargin=*]
  \item Do another full timed practice exam
  \item Identify weak topics
\end{itemize}

\textbf{Afternoon (2--3 hours):}
\begin{itemize}[leftmargin=*]
  \item Focus on problems you got wrong
  \item Review key formulas
  \item Create/refine your cheat sheet
\end{itemize}

\textbf{Evening:}
\begin{itemize}[leftmargin=*]
  \item Light review, no new material
  \item Get good sleep before exam!
\end{itemize}

\subsection{Emergency 2-Day Crash Plan}

If you only have 2 days:

\subsubsection{Day 1: Poisson + Random Walks ONLY (10--12 hours)}
\begin{itemize}[leftmargin=*]
  \item These are 100\% guaranteed to appear
  \item Master thinning and first-step analysis
  \item Do F2025 Q2, Q3 and S2025 Q1, Q3
\end{itemize}

\subsubsection{Day 2: Joint Densities + Full Practice (10--12 hours)}
\begin{itemize}[leftmargin=*]
  \item Learn conditional density basics
  \item Do 2 full timed practice exams
  \item Focus on understanding solutions
\end{itemize}

\textbf{Note:} Skip Brownian Motion if pressed for time. Focus on the big 3!

\subsection{How to Learn Each Topic Effectively}

\subsubsection{1. Read Slides First}
\begin{itemize}[leftmargin=*]
  \item Don't rush, understand each example
  \item Write down key formulas in your own notes
  \item Try to explain concepts in your own words
\end{itemize}

\subsubsection{2. Practice Problems Immediately}
\begin{itemize}[leftmargin=*]
  \item Don't just read solutions!
  \item Attempt problem yourself first (even if you struggle)
  \item Then check solution and understand each step
\end{itemize}

\subsubsection{3. Pattern Recognition}
\begin{itemize}[leftmargin=*]
  \item Notice: ``When I see X, I should use technique Y''
  \item Example: ``When I see `types of customers' $\rightarrow$ Thinning!''
  \item Example: ``When I see `hitting probability' $\rightarrow$ First-step analysis!''
\end{itemize}

\subsubsection{4. Active Recall}
\begin{itemize}[leftmargin=*]
  \item After learning a topic, close your notes
  \item Try to write key formulas from memory
  \item Explain the technique out loud to yourself
\end{itemize}

\subsubsection{5. Interleaved Practice}
\begin{itemize}[leftmargin=*]
  \item Don't do 10 Poisson problems in a row
  \item Mix topics: Poisson $\rightarrow$ Random Walk $\rightarrow$ Joint Density
  \item Helps you recognize which technique to use
\end{itemize}

\newpage

\section{Quick Reference: Which Technique When?}

\begin{table}[h]
\centering
\caption{Problem Type Recognition Guide}
\begin{tabularx}{\textwidth}{lX}
\toprule
\textbf{When you see...} & \textbf{Use this technique} \\
\midrule
``Poisson process'' + ``types/categories'' & \textbf{Thinning} \\
``random walk'' + ``probability of hitting'' & \textbf{First-Step Analysis (hitting prob)} \\
``random walk'' + ``expected time'' & \textbf{First-Step Analysis (expected time)} \\
``given that $N(t) = n$'' & \textbf{Conditional Expectation} with Poisson \\
``exponential'' + ``minimum'' or ``first to occur'' & \textbf{Min of Exponentials} \\
``joint density'' or ``$f(x,y)$'' & Set up \textbf{Joint PDF}, check range \\
``given $X = x$, find $E[Y]$'' & \textbf{Conditional Expectation/Conditional Density} \\
``Brownian motion'' or ``$B(t)$'' & Use \textbf{BM properties} ($E[B(t)] = 0$, Var $= t$) \\
\bottomrule
\end{tabularx}
\end{table}

\section{Practice Problem Recommendations}

\subsection{Must-Do Problems (Do these first!)}

\begin{enumerate}[leftmargin=*]
  \item \textbf{F2025 Q3} --- Poisson thinning + conditional expectation
  \item \textbf{F2025 Q2} --- Random walk with modified probabilities
  \item \textbf{S2025 Q1} --- Biased random walk first-step analysis
  \item \textbf{S2025 Q3} --- Competing Poisson processes
  \item \textbf{2022F Q1} --- Poisson with exponential threshold
  \item \textbf{2022S Q1} --- Poisson + exponential conditional
\end{enumerate}

\subsection{Good Practice Problems}

\begin{enumerate}[leftmargin=*,start=7]
  \item \textbf{F2025 Q1} --- Joint density $Y = UX$
  \item \textbf{S2025 Q2} --- Joint density with exponentials
  \item \textbf{2021 Q2} --- Poisson two-stage thinning
  \item \textbf{2019 Q3} --- Biased random walk
\end{enumerate}

\subsection{Advanced/Optional}

\begin{enumerate}[leftmargin=*,start=11]
  \item \textbf{2022S Q3} --- Brownian motion (stock prices)
  \item \textbf{2021 Q1} --- Two Brownian motions
  \item \textbf{2017 Q2} --- Hazard functions
\end{enumerate}

\newpage

\section{Exam Day Tips}

\begin{enumerate}[leftmargin=*]
  \item \textbf{Read the problem carefully} --- identify which technique to use
  \item \textbf{Write down what you know} --- given distributions, parameters, etc.
  \item \textbf{Set up the problem} --- define random variables clearly
  \item \textbf{Show your work} --- partial credit is awarded!
  \item \textbf{Check boundary conditions} --- especially for random walks
  \item \textbf{Simplify your answer} --- reduce fractions, combine terms
  \item \textbf{Sanity check} --- does your answer make intuitive sense?
\end{enumerate}

\section{Knowledge Check: Are You Ready?}

Before the exam, you should be able to:

\subsection{Poisson Processes}
\begin{itemize}[leftmargin=*]
  \item[$\square$] Explain what thinning means and when to use it
  \item[$\square$] Set up $P(\text{A before B})$ for competing processes
  \item[$\square$] Use $E[T \mid N=n]$ correctly
\end{itemize}

\subsection{Random Walks}
\begin{itemize}[leftmargin=*]
  \item[$\square$] Write first-step analysis formula from memory
  \item[$\square$] Remember to include ``+1'' for expected time formula
  \item[$\square$] Identify and write correct boundary conditions
  \item[$\square$] Solve a system of 2--3 linear equations
\end{itemize}

\subsection{Joint/Conditional Densities}
\begin{itemize}[leftmargin=*]
  \item[$\square$] Build joint PDF from conditional + marginal
  \item[$\square$] Specify correct range for density
  \item[$\square$] Use law of total expectation $E[Y] = E[E[Y|X]]$
  \item[$\square$] Find marginal from joint PDF
\end{itemize}

\subsection{Exponential}
\begin{itemize}[leftmargin=*]
  \item[$\square$] Know memoryless property
  \item[$\square$] Set up min of exponentials
  \item[$\square$] Remember $E[T] = 1/\lambda$ (not $\lambda$!)
\end{itemize}

\newpage

\section{Final Motivation}

\begin{tipbox}
\textbf{Remember:}
\begin{itemize}[leftmargin=*]
  \item You have a clear roadmap now!
  \item Topics are well-defined and learnable
  \item Focus on understanding, not memorizing
  \item Past students have done this successfully
  \item You can learn a lot in a focused few days
\end{itemize}
\end{tipbox}

\subsection{Strategy for Success}

\begin{enumerate}[leftmargin=*]
  \item Don't try to learn everything perfectly
  \item Focus on high-frequency topics first
  \item Understand the \textbf{TECHNIQUES}, not just formulas
  \item Practice with actual exam problems
  \item Learn from mistakes
\end{enumerate}

\subsection{You've Got This!}

The fact that you're being proactive and strategic shows you have what it takes to succeed.

\vspace{1cm}

\begin{center}
\large\textbf{Now stop reading and start learning!}

\vspace{0.5cm}

Good luck! You can do this!
\end{center}

\end{document}
